\section{第三人撤销之诉}
\begin{enumerate}
  \item 是一个新的\zd{形成之诉}。
  \item 适用\zd{一审普通程序}。
  \item 由作出该\zd{生效判决/终审判决}的法院管辖。
  \item 不予受理的情形。
        \begin{enumerate}
          \item 非讼案件。
          \item 解除婚姻关系涉及\zd{身份关系}的。保护到当事人的再婚权。
          \item 没有参加登记的权利人对代表人诉讼的生效裁判。因为可以另行起诉。
        \end{enumerate}
  \item 开庭审理，不适用简易程序，是一审判决。
  \item \zd{第三人是原告，原审的原告、被告、有独三、担责无独三作被告，不担责无独三作第三人}。
  \item 注意题目的干扰信息，不用看第三人的主张是什么。
  \item 中止执行。
        \begin{enumerate}
          \item 讲理：提\zd{案外人对执行标的的异议}，成立中止，不成立驳回 \yjt 继续三撤呗。
          \item 给钱：提供担保。
          \item 既不讲理又不给钱。你这种人又没有钱又不要脸，就没人理你。
        \end{enumerate}
  \item 怎么判。两个请求：\yi 判错了 \yjt 希望你们撤销 \er 房子是我的 \yjt 应该改判给我。
        \begin{enumerate}
          \item 是你的 \yjt \zd{撤销 + 改判}原判。
          \item 不是你的 \yjt \zd{撤销}原判错误部分。
          \item 完全不成立 \yjt 驳回诉讼请求。
        \end{enumerate}
  \item \zd{与再审的关系}。能够推翻一份生效裁判的只有这两个制度。
        \begin{enumerate}
          \item 绝对不能同时进行，\zd{互斥存在}。目的一样，功能一样。
          \item 启动主体不一样。\zd{再审：原审当事人 + 本应作原神共同原告、被告的}。
          \item 在三撤审理过程中，法院裁定再审。原则上\zd{再审优先}，三撤并入再审。
                例外：原审当事人恶意串通，三撤优先，中止再审。
          \item 怎么并：
                \begin{enumerate}
                  \item 适用一审的再审并进去一并审一并判。
                  \item 适用二审的再审并进去只准调不准判调不成撤销原判发回重审。
                \end{enumerate}
        \end{enumerate}
\end{enumerate}
